\section{不是五个朝代的尾声：九〇七至九六〇年的并行选择}
\label{sec:five-dynasties-curated-chronicle}

907年朱温废唐建后梁。\eventanchor{LT-907-080}{后梁建国（907）}这条中原年序立刻遇到河东李克用集团、吴、吴越、前蜀及北方契丹的并行行动；它不能替代南方和边地的时间。朝廷名号、军镇兵源、盐粮与水路城市从此在多个中心被重新组合。两《五代史》保存中原政权的连续性，却不足以让我们把吴、闽、南汉、楚、荆南、前后蜀、北汉写成背景。

同年，王建在成都称帝，杨渥拒受后梁册命，钱镠则接受吴越王的册封。\eventanchor{LT-907-082}{前蜀建立（907）}\eventanchor{LT-907-084}{吴继续奉唐年号（907）}\eventanchor{LT-907-085}{吴越受封（907）}这三种选择都不是对“正统”的简单表态：成都要把四川的军政资源组织为前蜀，淮南要处理杨氏、徐氏与唐朝名义的关系，杭州则以册封维系沿海和两浙的区域位置。916年阿保机称帝、917年刘龑在番禺称帝，又使草原与岭南各自取得新的年号和制度中心。\eventanchor{LT-916-089}{契丹建国（916）}\eventanchor{LT-917-090}{岭南称帝（917）}

923年李存勖建后唐、灭后梁；936年石敬瑭借契丹军力建后晋。\eventanchor{LT-923-092}{后唐灭后梁（923）}\eventanchor{LT-936-099}{后晋建国（936）}后者把北方政权继承和契丹的战略选择系在一起。传统叙事常将燕云问题压缩成一次“割让”的道德判决；更能看见的，是边地城镇、骑兵通道、契丹援军与新政权合法性在危机中彼此交换。具体条约条件、军数和地方反应，仍要保留材料的空白。

\begin{sourcenote}{十国的名号不等于全境控制}
前蜀、吴、吴越、契丹和南汉的建国年序可由两《五代史》、各国国史及《辽史》互校；这些书的保存密度和立场并不相同。册封、受禅或称帝能说明政治语言和中心位置，不能证明成都、杭州、番禺以外的州县已按同一方式纳税、服役或认同新政权。
\end{sourcenote}

\begin{event}{\eventref{LT-907-084}{杨渥续用唐号与淮南的选择}}{天祐四年（九〇七年）}{吴与后梁、吴越等政权 · 扬州、淮南水网及江南城市}{朱温建梁和杨氏续唐名号可核；各州服从、财赋规模和政治意图不宜由国号推断}
朱温在开封废唐建梁时，淮南的杨渥没有把新的中原国号当作唯一可用的政治语言。扬州周围的水网、盐粮、军镇和城市市场，使维持唐号既关乎合法性，也关乎如何同部将、地方官和往来商人维持可以预期的关系。这个选择并未把吴变成脱离北方的孤岛：后梁、吴越、前蜀和契丹的行动继续改变道路、使节和物资的去处。它的直接后果是中原更替不能独占年序；中期后果是江南诸政权须以不同的称号、税源和水路组织自己，而不是被动等待北方统一。

《旧五代史》《新五代史》《资治通鉴》可确认后梁建立、杨氏政权和吴越受封等基本次序，地方碑刻、墓志和钱币能补见若干城市与制度活动。材料的保存以中原为重，不能将“续用唐号”解释为淮南所有人共享的忠唐情感，亦不能由此量出财政或人口。读者应把九〇七年的名号选择放回扬州和江南的水陆条件中：它是一项维持区域关系的行动，而不是五代十国导航表上的一行标签。
\end{event}

\begin{causalchain}{北方继承与南方经营为何不能分开读}
中原政权需要军粮与边防支撑，契丹需要南下的伙伴与边地通道；这是使能条件。继承危机是触发，援军与城镇控制是直接后果。与此同时，江淮、两浙、岭南和四川的政权以运河、海路、盐茶、稻作和本地军队维持不同节奏；不能把其制度成绩或负担简单归为中原乱世的副产品。
\end{causalchain}

947年契丹入汴，随后后汉建立；951年郭威建后周。\eventanchor{LT-947-106}{契丹入汴（947）}\eventanchor{LT-951-109}{后周建立（951）}这些更替显示，开封控制并不足以立即取得各地的仓储、道路和忠诚。后周的整军、铸钱与对淮南的行动，必须放在吴越、南唐、北汉和契丹的选择旁边阅读，而不是预先当作宋代统一的铺垫。

南唐在937年由徐知诰受禅而立，945年攻取闽地建州；两次转换都须放在江淮、福建的地方军政与水陆通道中理解。\eventanchor{LT-937-100}{南唐建立（937）}\eventanchor{LT-945-104}{南唐取闽（945）}951年刘崇在太原建北汉并求援于辽，954年高平之战把后周、北汉和辽的北方竞争推到战场上。\eventanchor{LT-951-110}{北汉建立（951）}\eventanchor{LT-954-111}{高平之战（954）}956--957年后周南征，南唐割让江北诸州；条约并不是地方接收已完成的证明。\eventanchor{LT-956-113}{后周入淮南（956）}\eventanchor{LT-957-114}{南唐割江北（957）}

\begin{event}{\eventref{LT-937-100}{徐知诰受禅与南唐的江淮接手}}{升元元年（九三七年）}{南唐与江宁、扬州及江淮水陆网络}{受禅、改名和建国年序可核；州县服从、税源规模与地方意图不能由新国号推断}
徐知诰从吴政权接手名义时，江宁的政治中心必须同扬州、江淮的军镇、盐粮和水路城市维持可用的联系。改国号给任命、文书与外交提供了新的名义，却没有使原有将校、地方官和商人同时换成一种政治关系；南唐能调度什么，仍取决于河道、仓储和区域人际网络是否接得住。后来向闽地推进以及与后周在淮南的冲突，都把这项接手的限度逐一暴露出来。

《旧五代史》《新五代史》与《资治通鉴》能够交叉确定政权转换的大致顺序，南唐相关国别材料、墓志、钱币和城市遗存可补见若干地方活动。国史散佚和中原材料的重心意味着，我们不能据受禅文辞判定江淮所有州县的认同，也不能把局部遗存换算为南唐财政或人口的总量。
\end{event}

960年赵匡胤取代后周。\eventanchor{LT-960-118}{陈桥与宋代建立（960）}这结束的是后周在开封的政权，不是同年消失的十国世界。北汉仍在，南唐、吴越、南汉、闽地遗绪与蜀地秩序各有延续；统一将在日后的战争、协商和财政接管中逐步发生。

\begin{sourcenote}{五代十国的材料为何特别不均衡}
《旧五代史》《新五代史》《资治通鉴》及《考异》维持中原年序；地方碑刻、墓志、钱币、城址和残存文书使南方与边地的一部分实践可见。国史散佚和后来的宋代视角意味着：对十国的财政、人口与社会，常只能作有限判断，不能把沉默写成不存在。
\end{sourcenote}

本章的路线可回看\eventref{LT-880-062}{黄巢入长安（880）}怎样加速区域军府的形成，也可从\eventref{ST-605-TONGJI-CANAL}{通济渠（605）}重新审视水路、粮运与国家能力的长线。这样的往返不是把事件表串成链，而是检验不同政权何时共享一条道路、何时又各自承担其成本。

\begin{event}{\eventanchor{LT-924-093}{高季兴受封南平王：江陵不是后唐的一个省略号}}{同光二年（九二四年）}{荆南与后唐 · 江陵、长江中游与荆襄水陆要冲}{高季兴受封与江陵政权的存在可核；实际辖境、税源和周边州县服从程度不能由王号断定}
后唐册封高季兴为南平王时，江陵的价值不在一枚王印，而在它守着长江中游的船路、荆襄之间的陆道和可以转运粮盐的城市节点。高季兴既需要借后唐的名义使本地将校、商旅与邻近势力知道自己有可用的政治位置，也不能把江陵的仓储、兵员和水路完全交给北方朝廷调度。对后唐而言，承认荆南是一种降低即时军事成本、维持通道可用的选择；对江陵的军民和地方官而言，册封之后谁来征发、谁负责城防和谁能保障航运，才决定新关系的重量。这个小政权并不是五代中原更替之间的停顿，而是各方必须经过或绕开的中介。

《旧五代史·高季兴传》《新五代史·南平世家》与《十国春秋》能确认受封及其基本次序，但后者成书较晚，三者都不能把王号直接换算成稳定的州县控制。江陵碑志、墓葬、钱币和城市遗存只能补见局部活动，也不足以给出完整财政账。受封的直接效果是后唐与荆南获得一套可协商的名义；中期效果却是江陵得以在后唐、南唐、吴越等势力之间保存自己的调度空间。它的代价是每一次北方战争、江上封锁或邻国使节变化，都可能使这套依赖道路的平衡重新被议价。
\end{event}

\begin{event}{\eventanchor{LT-938-BACH-DANG}{白藤江的潮水：南汉水军为何未能重建交州关系}}{天福三年（九三八年）}{交州与南汉 · 白藤江水口、红河三角洲与岭南海路}{吴权击败南汉水军、南汉未恢复对交州的主干可核；战术细节、兵数和潮汐利用方式不能作定论}
白藤江把岭南与交州之间的海路压缩到一个会随潮汐、河口地形和地方引导而改变的水面。南汉若要由广州方向恢复对交州的影响，军舰不仅得抵达河口，还得维持补给、取得可靠的水路信息，并让登陆后的部队能同地方城镇发生可持续的联系。吴权一方能够利用熟悉河道的条件抵抗，使南汉的远征没有变成稳定接收；战后交州的政治选择也不能被理解为一场水战自动创造的共同意志。船工、沿岸居民和地方首领面对的是港口、税源与保护者的重新安排，而不仅是两国史书里的胜负。

《大越史记全书》以后来建国叙事保存此战，《新五代史·南汉世家》和《资治通鉴》则从南汉及中原编年位置简略记录，两类材料可以相互约束年序，不能拼出一份逐刻战报。木桩、河口地形和考古研究能讨论战场条件，却不允许把舰船数、伤亡或指挥口令填成确定事实。白藤江的直接后果是南汉未能重建对交州的军事控制；更长的影响则是岭南、交州和海上商路继续在不同政治中心之间调整。这使十国史不再只是开封更替的余波，而是一段由河口、船队和地方协作决定的并行年序。
\end{event}
